\section*{Problem 1}

在 $\mathbb{R}^3$ 中，设子空间 $W_1$ 由方程 $x+y-z=0$ 确定，子空间 $W_2$ 由方程 $x-y=0$ 确定.请分别计算子空间 $W_1$、$W_2$ 以及交空间 $W_1 \cap W_2$ 的维数，并利用维数公式求和空间 $W_1 + W_2$ 的维数.

\section*{Problem 2}

设 $\mathbb{R}^4$ 中的两个子空间分别为
\[
\begin{aligned}
    U &= \Span\left\{
    \left(1, 0, 0, 0\right)^\top,
    \left(0, 1, 0, 0\right)^\top
    \right\}, \\
    V &= \Span\left\{
    \left(1, 1, 1, 1\right)^\top,
    \left(0, 0, 1, 1\right)^\top
    \right\}.
\end{aligned}
\]
请通过计算维数 $\dim \left(U + V\right)$，来判断 $U + V$ 是否为直和（即判断是否满足 $U \oplus V$）.

\section*{Problem 3}

已知 $n$ 阶方阵空间 $V = \mathbb{R}^{n \times n}$ 可以分解为对称矩阵子空间 $W_1$ 与反对称矩阵子空间 $W_2$ 的直和（即 $V = W_1 \oplus W_2$）.给定矩阵
\[
A =
\begin{bmatrix}
1 & 2 & 4 \\
4 & 5 & 6 \\
2 & 8 & 9
\end{bmatrix},
\]
请利用直和分解的性质，求出唯一确定的对称矩阵 $S \in W_1$ 和反对称矩阵 $K \in W_2$，使得 $A = S + K$.
